\usepackage[paperwidth=178mm,paperheight=254mm,textwidth=117mm,hcentering,
  top=20mm,bottom=27mm,headheight=14pt,headsep=4mm,footskip=13mm]{geometry}
\usepackage{amsmath,cancel}
\usepackage{fontspec,xeCJK,unicode-math}
\setmainfont{TeX Gyre Termes}
\setsansfont{TeX Gyre Heros}
\setmonofont{TeX Gyre Cursor}
\setmathfont{TeX Gyre Termes Math}
\setCJKmainfont{FandolSong-Regular.otf}[
  BoldFont=FandolSong-Bold.otf,ItalicFont=FandolKai-Regular.otf]
\setCJKsansfont{FandolHei-Regular.otf}[BoldFont=FandolHei-Bold.otf]
\setCJKmonofont{FandolFang-Regular.otf}
\xeCJKsetup{PunctStyle=kaiming}
\usepackage{graphicx,booktabs,longtable,tabularx,array,multicol,pdflscape}
\usepackage[table]{xcolor}
\usepackage{titlesec,enumitem,fancyhdr,listings,caption}
\usepackage[most]{tcolorbox}
\usetikzlibrary{arrows.meta,positioning,calc,shapes.geometric,shapes.arrows,fit,backgrounds}
\usepackage{xurl}
\usepackage[unicode,hidelinks,bookmarksnumbered,bookmarksdepth=4]{hyperref}
\usepackage{bookmark}
\hypersetup{pdftitle={大语言模型深度解析——贯通理论与实践},
  pdfauthor={Uday Kamath, Kevin Keenan, Garrett Somers, Sarah Sorenson}}

\renewcommand\contentsname{目录}
\renewcommand\chaptername{第}
\renewcommand\appendixname{附录}
\renewcommand\figurename{图}
\renewcommand\tablename{表}
\renewcommand\lstlistingname{清单}
\setcounter{secnumdepth}{4}
\setcounter{tocdepth}{2}

\titleformat{\chapter}[display]{\sffamily\bfseries\Large}{\ifinappendix 附录~\thechapter\else 第~\thechapter~章\fi}{6pt}{\Huge}
\titleformat{name=\chapter,numberless}[block]{\sffamily\bfseries\Huge}{}{0pt}{}
\titlespacing*{\chapter}{0pt}{0pt}{32pt}
\titleformat{\section}{\sffamily\bfseries\Large}{\thesection}{.6em}{}
\titleformat{\subsection}{\sffamily\bfseries\large}{\thesubsection}{.6em}{}
\titleformat{\subsubsection}{\sffamily\bfseries\normalsize}{\thesubsubsection}{.6em}{}
\titlespacing*{\section}{0pt}{20pt plus 3pt}{9pt}
\titlespacing*{\subsection}{0pt}{16pt plus 3pt}{8pt}
\titlespacing*{\subsubsection}{0pt}{13pt plus 3pt}{7pt}
\titleformat{\paragraph}[hang]{\sffamily\bfseries\normalsize}{\theparagraph}{.6em}{}
\titlespacing*{\paragraph}{0pt}{10pt plus 2pt}{5pt}
\newif\ifinappendix
\AddToHook{cmd/appendix/after}{\inappendixtrue}
\renewcommand\normalsize{\fontsize{11pt}{16.5pt}\selectfont}
\AtBeginDocument{\normalsize}
\setlength{\parindent}{2em}
\setlength{\parskip}{0pt}
\setlength{\emergencystretch}{1.5em}
\setlength{\columnsep}{16pt}
\raggedbottom
\renewcommand\floatpagefraction{.72}
\clubpenalty=8000
\widowpenalty=8000
\setlist{leftmargin=2em,itemsep=2pt,topsep=6pt,parsep=0pt}

\pagestyle{fancy}
\fancyhf{}
\fancyhead[LE,RO]{\small\thepage}
\fancyhead[RE]{\small\nouppercase{\leftmark}}
\fancyhead[LO]{\small\nouppercase{\rightmark}}
\renewcommand\headrulewidth{0pt}
\renewcommand\chaptermark[1]{\markboth{\thechapter\enspace #1}{\thechapter\enspace #1}}
\renewcommand\sectionmark[1]{\markright{\thesection\enspace #1}}
\fancypagestyle{plain}{\fancyhf{}\fancyfoot[LE,RO]{\small\thepage}}

\newcommand{\frontchapter}[1]{%
  \chapter*{#1}\pdfbookmark[0]{#1}{front.\thepage}\markboth{#1}{#1}}
\newcommand{\fronttocchapter}[1]{%
  \chapter*{#1}\addcontentsline{toc}{chapter}{#1}\markboth{#1}{#1}}
\newcommand{\frontsection}[1]{%
  \section*{#1}\pdfbookmark[1]{#1}{frontsec.\thepage.\number\value{section}}%
  \stepcounter{section}}
\newcommand{\chapterabstract}[1]{\noindent\textbf{摘要}\quad #1\par\medskip}
\newcommand{\chaptercredit}[1]{%
  \begingroup\renewcommand\thefootnote{}\footnotetext[0]{%
    \textcopyright\ 作者，2024 年独家许可予 Springer Nature Switzerland AG。%
    U. Kamath 等，\textit{Large Language Models: A Deep Dive}。
    \if\relax\detokenize{#1}\relax
      \url{https://doi.org/10.1007/978-3-031-65647-7}%
    \else
      \url{https://doi.org/10.1007/978-3-031-65647-7_#1}%
    \fi}%
  \endgroup}
\newcommand{\reviewer}[1]{\par\smallskip\hfill\null\hfill
  \parbox{.9\linewidth}{\raggedleft\small #1}\par\medskip}

\newtcolorbox{keypoint}{enhanced,breakable,colback=black!14,colframe=black!14,
  arc=0pt,boxrule=0pt,left=9pt,right=9pt,top=6pt,bottom=6pt,
  before skip=8pt,after skip=8pt}
\newtcolorbox{practicaltip}{enhanced,breakable,colback=white,colframe=black!34,
  colbacktitle=black!30,coltitle=black,arc=0pt,boxrule=.35pt,
  title={实用提示},fonttitle=\sffamily\bfseries,
  left=6pt,right=6pt,top=5pt,bottom=5pt,before skip=8pt,after skip=8pt}
\lstdefinestyle{booklisting}{basicstyle=\ttfamily\fontsize{8.6pt}{11.5pt}\selectfont,columns=fullflexible,
  keepspaces=true,showstringspaces=false,commentstyle=\color{green!40!black},breaklines=true,frame=single,
  rulecolor=\color{black!45},backgroundcolor=\color{black!4},captionpos=b,
  aboveskip=8pt,belowskip=8pt}
\lstset{style=booklisting}
\captionsetup{font=small,labelfont=bf,labelsep=quad}
\newcommand{\sourcefigure}[3]{%
  \begin{figure}[!htbp]
  \centering
  \includegraphics[width=\linewidth,height=.65\textheight,keepaspectratio]{assets/#1}%
  \caption{#2}\label{#3}%
  \end{figure}}
\newenvironment{chapterreferences}{%
  \section*{参考文献}\phantomsection\addcontentsline{toc}{section}{参考文献}%
  \markright{参考文献}\small\raggedright\setlength{\parindent}{0pt}%
  \setlength{\parskip}{3pt}\everypar{\hangindent=1.4em\hangafter=1}%
}{\par}
\newcommand{\transnote}[1]{\par\smallskip\noindent{\small\textbf{译注}\quad #1\par}\smallskip}
\newcommand{\abb}[3]{#1 & #2 & #3\\}

\newcommand{\refentry}[2]{\par\noindent\phantomsection\label{bib:#1}#2\par}
\DeclareRobustCommand{\refcite}[2]{\hyperref[bib:#1]{#2}}
